% s-345-p345-m002; original PDF 345, 346, 347
\Aphorism{s-345-p345-m002}{398}
Die Abweichungen der Planetenbahnen von der reinen Kegelschnittform nennt der Astronom „Störungen“: welche naive Anmassung! Was eine Störung seiner rechnerischen Bequemlichkeit bedeutet, empfindet er unbewusst auch als Störung des eigentlichen gesetzmässigen Naturverlaufs, obwohl es nichts anderes als dieser Verlauf selber ist, nur verfolgt bis in seine letzten „unangenehmsten“ Consequenzen. Früher war man, der grösseren Unwissenheit entsprechend, noch unbescheidener; da stellte bereits der Kegelschnitt, dem Kreise gegenüber, eine „Störung“, eine Complication ersten Ranges vor, und die Leistung des Planeten, der in einer Ellipse um die Sonne wandelt, galt noch Geistern wie Bruno und Vanini für undenkbar ohne irgendwelche planetarische Intelligenz. Im Grunde urtheilen wir Alle heute noch über Einfachheit und Complication ebenso anmassend und anthropomorphistisch; was uns in der gedanklichen Nachbildung am wenigsten Kopfzerbrechen macht, das hätte, meinen wir, auch der Natur in der factischen Verwirklichung am wenigsten Mühe verursacht, das also hätte sie \Sperr{eigentlich} gewollt und auch erreicht, wenn ihr nicht die „Störungen“ dazwischen gekommen wären. Wir stecken voll scholastischen und aristotelischen Aberglaubens; der Kreis ist \Sperr{eigentlich} die Bahn der Himmelskörper und der goldene Schnitt das \Sperr{eigentliche} Bildungsgesetz in der Natur. Selbst unser moderner Entwicklungsbegriff gehört zu diesem Aberglauben; die Natur als Elevin, die sich zur Meisterin durchstümpert, die \Sperr{zuerst} Moneren und Zellen und \Sperr{dann} in steigender Differenzirung höhere Formen hervorbringt — was in aller Welt berechtigt uns zu dieser allzumenschlichen Metapher! Ich meine, die Natur sprudelt alles heraus, mit gleicher Fertigkeit und Spielerei, das Dümmste ebenso gern wie das Gescheidteste; wir haben die Natur nicht erklärt, sondern uns nur menschlich näher gebracht, wenn wir das Complicirte auf Einfacheres zurückführen. Für unsere noch so unbeholfene Analysis müssen sogar die complicirten Erscheinungen die Regel sein und die einfachen nur Ausnahme. Nicht die „Störungen“ sind ein Phänomen, sondern die Kleinheit der Störungen; nicht dass der Wasserfall keine stationäre Strömung ist, verdient unser Staunen, sondern dass er noch so \Sperr{nahezu} stationär ist.
\par

